\chapteropen{04}

\begin{cele}
  \item prowadzić dwie listy: nieograniczoną listę „do przeczytania” i kolejkę z trzema pozycjami, które naprawdę czekają na twoją uwagę,
  \item stosować regułę porzucania opartą na liczbie stron, a nie na nastroju,
  \item ułożyć plan dwunastu pozycji, który przeplata lektury lekkie, średnie i trudne,
  \item policzyć, ile książek krótkich i typowych mieści się w roku, i wiedzieć, kiedy ta liczba coś mówi, a kiedy nic,
  \item rozpoznać, kiedy zakup książki zwiększa stosik, a nie czytanie.
\end{cele}

Wybór książki jest decyzją, którą podejmujesz kilka razy w tygodniu, często bez zastanowienia: w księgarni, w bibliotece, po polecenie od znajomego, na liście w serwisie z ofertami. Każda taka decyzja zabiera część czasu, który policzyliśmy w rozdziale o arytmetyce, i każda ma koszt, którego zwykle nie widać. Ten rozdział nie podaje listy tytułów. Podaje procedury, dzięki którym wybór zostaje w twoich rękach, a nie w rękach tego, co ostatnio pojawiło się na ekranie.

\section{Dwie listy: do przeczytania i kolejka}

Potrzebujesz dwóch list, bo robią dwie różne rzeczy. Pierwsza to \glqq do przeczytania\grqq{}: wszystko, co kiedyś chcesz przeczytać, bez limitu i bez terminów. Ta lista służy do zbierania. Jej zadaniem jest to, żeby pomysł nie przepadł, a nie to, żeby czekał na dzień, w którym będziesz gotowy. Jeśli zapisujesz każdą wzmiankę o książce, lista rośnie szybko, i to jest w porządku.

Druga lista to kolejka, i ma trzy pozycje. Nie więcej. W kolejce są wyłącznie te książki, które w najbliższych tygodniach zamierzasz otworzyć. Gdy kończysz jedną z nich, z listy do przeczytania wchodzi następna, a wybiera ją ty, a nie algorytm sklepu ani kolejność, w jakiej coś wpadło ci w oko.

Limit trzech nie jest magiczny. Wynika z prostego rachunku. Jeśli czytasz dziesięć minut dziennie, przy typowej długości \val{typical.words} słów przerobisz średnio \val{books.10.250} książki w roku. Kolejka z dziesięcioma pozycjami oznacza, że większość z nich przez cały rok nie zostanie otwarta, bo po prostu nie ma czasu, żeby je zacząć. Trzy pozycje to tyle, ile da się realnie prowadzić przy jednym czytaniu dziennie: jedna, którą czytasz teraz, i dwie, które czekają na swoją kolej.

Lista do przeczytania może być długa, kolejka nie może. Ta różnica jest najważniejszą regułą tego rozdziału, bo chroni przed najczęstszą pomyłką: traktowaniem wszystkiego, co chciałbyś przeczytać, jak czegoś, co już zaczynasz.

\section{Zasada pięćdziesięciu stron}

Potrzebujesz reguły porzucania, bo bez niej decyzję podejmuje nastrój. A nastrój zwykle podejmuje ją w złą stronę: po ciężkim dniu każda książka wydaje się zbyt ciężka, a po dobrym tygodniu zbyt płytka, i w obu przypadkach odkładasz ją bez sprawdzenia. Poniższa reguła jest liczbowa i dotyczy wyłącznie tego, kiedy sprawdzasz. Nie mówi, czy książka jest dobra.

\begin{zasada}[Próg sprawdzenia]
  Próg wynosi mniejszą z dwóch liczb: 50 stron albo 10\% długości książki. Gdy go przekroczysz, zadajesz jedno pytanie: czy chcę do tej książki wracać? Jeśli nie, porzucasz ją bez tłumaczenia się przed kimkolwiek, także przed sobą. Jeśli tak, czytasz dalej. Reguła nie zmusza cię do porzucania książki, którą chcesz czytać.
\end{zasada}

Dlaczego mniejsza z dwóch liczb? Dla książek krótszych niż 500 stron decyduje dziesięć procent, bo ograniczenie do 50 stron byłoby zbyt późnym sprawdzeniem: zainwestowałbyś w krótką książkę znacznie więcej, zanim zapytasz o sens. Dla książek dłuższych limit 50 stron chroni przed tym, żeby przeczytać sto stron, zanim cokolwiek się wyjaśni. Liczba stron jest tu tylko przybliżeniem długości: w ebooku liczy się procent postępu, a nie numer strony.

\begin{przyklad}[Trzy progi]
  Książka na 320 stron: 10\% to 32 strony. Mniejsza z liczb 32 i 50 to 32, więc próg leży na stronie 32.

  Książka na 480 stron: 10\% to 48 stron. Mniejsza z liczb 48 i 50 to 48, więc próg leży na stronie 48.

  Książka na 800 stron: 10\% to 80 stron. Mniejsza z liczb 80 i 50 to 50, więc próg leży na stronie 50.
\end{przyklad}

Widać, że reguła zmienia się płynnie. Po przekroczeniu progu nie oceniasz całej książki, tylko decydujesz, czy warto zapłacić dalszym czasem za to, co już wiesz. Pytanie \glqq czy chcę wracać?\grqq{} zwykle daje odpowiedź szybciej niż pytanie \glqq czy to dobra książka?\grqq{}.

\begin{falle}[Skończę, bo zacząłem]
  Książkę trzeba dokończyć, bo już sporo w nią włożyłem. Włożone godziny nie wracają, niezależnie od tego, czy będziesz czytać dalej. Liczy się tylko to, ile zajmie reszta i czy wolisz ten czas na tę książkę, czy na pozycję z kolejki. Dokończenie z poczucia straty jest najdroższym sposobem na zmarnowanie czasu, który i tak został wydany.
\end{falle}

Porzucona książka nie jest porażką rocznego planu. W rozliczeniu, o którym mowa w rozdziale o formatach, nie wlicza się ona do przeczytanych, i dobrze, bo rozliczenie ma mierzyć ukończenia, a nie liczbę otwarć. Zwolnione miejsce w kolejce przechodzi do następnej książki z listy, i tyle.

\section{Dobór trudności: lekka, średnia, trudna}

Dobór książek to nie tylko kwestia tematu. Ta sama liczba godzin wygląda inaczej, gdy kolejne pozycje są ciężkie jedna po drugiej, a inaczej, gdy trudność się przeplata. Trzy kategorie definiuję sprawdzalnie, przez to, czego książka od ciebie wymaga:

\begin{itemize}
  \item \textbf{Lekka} to książka, którą czytasz zmęczony i którą po tygodniowej przerwie podejmujesz bez zgubienia wątku. Nie musisz wracać do poprzedniego rozdziału, żeby zrozumieć następny.
  \item \textbf{Średnia} wymaga pełnej uwagi, ale nie notatek. Czasem sprawdzisz pojęcie, a argument trzyma się w pamięci przez cały rozdział.
  \item \textbf{Trudna} nie da się czytać bez notatek, przerw na sprawdzenie i powrotów. Często czytasz ją dwa razy albo w kawałkach, a pierwsze podejście bywa tylko rozpoznaniem terenu.
\end{itemize}

Przeplatanie to reguła praktyczna, a nie wynik badań, i traktuj ją jako punkt wyjścia. Ciąg trudnych książek prowadzi zwykle do przerwy, bo każda kolejna zaczyna się od zmęczenia po poprzedniej. Ciąg lekkich nie zatrzymuje wprawdzie czytania, ale nie rozwija też tego, co już umiesz. Naprzemienność daje każdej trudnej książce jedną lżejszą pozycję przed sobą i jedną po sobie.

Poniżej plan na dwanaście pozycji, czyli na jedną książkę miesięcznie. Długość przypisana w planie jest planowaną wartością, nie wynikiem. Zwróć uwagę, że w tym planie trzy pozycje są krótkie, sześć typowych i trzy długie, więc średnia długość wynosi dokładnie typową, a rachunek z poprzedniego rozdziału daje dla tego planu \val{need.12} minuty dziennie.

\begin{center}
\begin{tabular}{@{}rllp{6.4cm}@{}}
\toprule
\textbf{Poz.} & \textbf{Trudność} & \textbf{Długość} & \textbf{Rola w planie} \\
\midrule
1  & lekka   & krótka  & rozgrzewka po przerwie \\
2  & średnia & typowa  & pierwsza pozycja z argumentem \\
3  & trudna  & typowa  & pierwsze wyzwanie, z notatkami \\
4  & lekka   & krótka  & odpoczynek po trudnej \\
5  & średnia & typowa  & argument bez notatek \\
6  & trudna  & długa   & najdłuższa pozycja, dwa podejścia \\
7  & lekka   & typowa  & lektura na wakacje lub na wyjazd \\
8  & średnia & typowa  & praca nad konkretnym pytaniem \\
9  & trudna  & typowa  & rozdział o metodzie, z ćwiczeniami \\
10 & lekka   & krótka  & przerwa przed końcem roku \\
11 & średnia & długa   & długa, ale czytelna narracja \\
12 & trudna  & długa   & koniec roku, powrót do pozycji 3 i 6 \\
\bottomrule
\end{tabular}
\end{center}

Ta tabela nie jest listą tytułów. Wpisujesz do niej własne książki, a kategorie trudności przydzielasz po pierwszych dwudziestu stronach, nie przed nimi. Pomyłka w przydziale jest zwykle widoczna od razu: lekka okazuje się trudna, więc przesuwasz ją do kolumny trudnych, i plan wraca do równowagi.

\section{Krótkie książki jako przyspieszacz}

Krótkie książki mają jedną zaletę, której nie da się przecenić, jeśli dopiero zaczynasz: zamykają się szybko. Ukończenie daje sygnał, że cel jest osiągalny, a sygnał taki bywa ważniejszy niż same słowa. Poniższa tabela pochodzi z danych w repozytorium, a nie z pamięci. Przelicza trzy rodzaje książek na słowa i godziny przy tempie 250 słów na minutę, czyli przy założeniu przyjętym w tej książce.

\generated{book-lengths}

Czytając tabelę, zapamiętaj jedną proporcję. Krótka książka ma 40\,000 słów, czyli połowę typowej. Pięć krótkich książek to 200\,000 słów, czyli dwa i pół typowej książki. Pięć ukończeń wygląda w statystyce lepiej niż jedno, ale za nimi stoi dwa i pół raza więcej czytania niż za jedną typową pozycją.

Przyspieszacz działa tylko wtedy, gdy nie zmienia celu. Gdy wybierasz krótkie książki wyłącznie po to, żeby liczba rosła, miara zaczyna sterować lekturą, a nie odwrotnie. Dlatego rozlicz dwie rzeczy razem: liczbę ukończonych książek i liczbę słów. Rok z dwudziestoma krótkimi pozycjami to osiemset tysięcy słów, rok z dwudziestoma typowymi to milion sześćset tysięcy. Oba mogą być dobre, ale nie są takie same.

Krótkie książki warto więc traktować jak rozgrzewkę i nagrodę przy trudnym tygodniu, a nie jak sposób na to, żeby wyglądać lepiej w podsumowaniu grudnia.

\section{Non-fiction i fikcja}

Format nie zmienia reguły porzucania, ale zmienia pytanie, które zadajesz po przekroczeniu progu. W literaturze faktu zwykle pytasz, czy argument jest użyteczny i czy zostaje w pamięci. Książka o strukturze rozdziałów pozwala czytać selektywnie: pierwsze kilkanaście stron, spis treści i wstęp mówią dużo o tym, czy warto dalej. Jeśli rozdział nie działa, możesz go przeskoczyć i wrócić do niego później, a książka nadal ma sens.

W fikcji pytasz, czy chcesz się dowiedzieć, co będzie dalej. To pytanie łatwo wyprzeć, bo pierwsze strony powieści często poświęcone są wprowadzeniu i nastrojowi, a nie akcji. Przerwa kosztuje w fikcji więcej: trzeba odtworzyć, kto jest kim, jaki jest stan sprawy i w jakim tonie prowadzona jest narracja. Dlatego po dłuższej przerwie wracaj do ostatniego rozdziału na chwilę, zanim zaczniesz nowy. To kilka minut, które oszczędzają wrażenia, że książka jest obca.

Oba formaty łączą się w planie z rozdziału o doborze trudności. Nie ma reguły, że fikcja jest lekka, a faktu trudny. Wiele powieści wymaga czytania uważnego, a wiele rozpraw popularnych czyta się płynnie. Kategoria zależy od tego, czego książka od ciebie wymaga, a nie od tego, jak została nazwana na okładce.

\section{Czytanie z polecenia i z rankingu}

Są dwa główne źródła pomysłów. Pierwsze to polecenie: ktoś, kogo znasz, mówi, że konkretna książka pomogła mu w konkretnej sprawie. Drugie to ranking: lista popularnych tytułów, zwykle ułożona według sprzedaży, ocen albo tego, co właśnie jest modne. Ranking mierzy, co wiele osób wybrało. Nie mierzy, czy ty potrzebujesz tego, co wybrały.

Polecenie ma zwykle więcej kontekstu, bo ktoś zna twoje sprawy albo przynajmniej wie, dlaczego sam sięgnął po książkę. To moja ocena, a nie wynik badania: polecenie od osoby, która zna twój cel, bywa bardziej użyteczne niż lista najczęściej kupowanych. Ranking jest jednak dobrym źródłem, gdy nie masz żadnego punktu odniesienia i chcesz zobaczyć, w jakiej rozmowie uczestniczysz.

Niezależnie od źródła zanim cokolwiek trafi do kolejki, odpowiedz na trzy pytania: kto poleca i z jakiego powodu, czego ja szukam w tej książce i dlaczego akurat teraz. Jeśli nie umiesz odpowiedzieć na drugie, książka trafia do listy do przeczytania, a nie do kolejki. Polecenie samo w sobie nie czyni książki ważniejszą od tej, która już na ciebie czeka.

\section{Stosik wstydu}

Stosik wstydu to książki, które masz, zamówiłeś albo zapisałeś, a nie przeczytałeś. Może leżeć na półce, w telefonie albo w koszyku. Problem nie polega na tym, że stosik istnieje, tylko na tym, że rośnie szybciej, niż czytasz, i stopniowo zaczyna rządzić decyzjami: zamiast wybrać książkę z kolejki, wyrzucasz sobie, że tyle już kupiłeś.

Rachunek jest prosty. Jeśli kupujesz dwie książki miesięcznie, to w roku wychodzi 24 zakupy. Przy dziesięciu minutach dziennie czytasz średnio \val{books.10.250} książki rocznie. Różnica między tymi liczbami rośnie co roku. Nie musisz kupować mniej, żeby to zmienić, wystarczy regułę wiązać z kolejką: nowa książka wchodzi do sklepu dopiero wtedy, gdy zwolni się miejsce w kolejce, czyli gdy kończysz poprzednią.

\begin{uwaga}[Zakup nie jest czytaniem]
  Kupienie książki daje przyjemność, a czasem poczucie, że już coś zrobiłeś. Nie zmienia jednak tego, ile godzin masz w tygodniu. Jeśli zakupy idą szybciej niż czytanie, zapisz w jednym miejscu liczbę książek w stosiku i sprawdzaj ją raz w miesiącu. Liczba ma być spadająca albo stała. Wzrost jest sygnałem do zmiany reguły, a nie do wyrzutów sumienia.
\end{uwaga}

Stosik nie jest karą. To miejsce, w którym widać, czy plan jest realistyczny, czy po prostu zbyt optymistyczny. Gdy liczba rośnie tygodniami, zmniejsz kolejkę albo zwiększ czas czytania; nie zmniejszaj wstydu.

\begin{cwiczenia}
  \exitem{Masz książkę na 320 stron. Po której stronie stosujesz regułę sprawdzenia?}{Po stronie 32. Dziesięć procent z 320 to 32, a mniejsza z liczb 32 i 50 to 32.}
  \exitem{Ile minut dziennie potrzeba, by w roku przeczytać 24 typowe książki przy 250 słowach na minutę?}{Około \val{need.24} minut dziennie. Rachunek: $24 \cdot 80\,000 = 1\,920\,000$ słów, $1\,920\,000 / 250 = 7\,680$ minut, $7\,680 / 365 \approx 21$ minut dziennie.}
  \exitem{W planie dwunastu pozycji, który przedstawiłem w tym rozdziale, które numery są książkami trudnymi?}{Pozycje 3, 6, 9 i 12. Każda trzecia pozycja w cyklu trzech jest trudna.}
  \exitem{Znajomy poleca ci książkę. Nie wiesz jeszcze, czego w niej szukasz. Gdzie ją wpisujesz?}{Do listy „do przeczytania”. Do kolejki trafia dopiero wtedy, gdy znasz odpowiedź na pytanie, czego szukasz, i gdy zwalnia się miejsce.}
  \exitem{Porzucasz po 50 stronach książkę na 800 stron. Czy to jest zgodne z regułą?}{Tak. Dziesięć procent z 800 to 80, a mniejsza z liczb 80 i 50 to 50. Próg leży na stronie 50, więc porzucenie po niej jest zgodne z regułą.}
  \exitem{Porównaj pięć krótkich książek z jedną typową. Ile słów i ile godzin przy 250 słowach na minutę daje każda z opcji?}{Pięć krótkich: 200\,000 słów, czyli 800 minut, około 13,3 godziny. Jedna typowa: 80\,000 słów, czyli 320 minut, około 5,3 godziny. Pięć ukończeń to więcej w statystyce, ale za nimi jest dwa i pół raza więcej słów niż za jedną typową.}
  \exitem{Rok z dwudziestoma krótkimi książkami i rok z dwudziestoma typowymi. Który wymaga więcej czasu i dlaczego liczba ukończeń nie wystarcza do porównania?}{Rok z typowymi wymaga więcej czasu: $20 \cdot 80\,000 = 1\,600\,000$ słów wobec $20 \cdot 40\,000 = 800\,000$. Liczba ukończeń jest ta sama, więc nie mówi, ile czytania za nią stoi. Trzeba rozliczać słowa razem z ukończeniami.}
\end{cwiczenia}

\zadanie{Zbuduj dwie listy: „do przeczytania” z co najmniej dwunastoma pozycjami i kolejkę z trzema. Każdej pozycji w kolejce przypisz trudność i długość. Przez cztery tygodnie czytaj pierwszą pozycję z kolejki i stosuj regułę progu. Miara sukcesu: po czterech tygodniach w kolejce są co najwyżej trzy pozycje, każda ma zapisany próg sprawdzenia, a stosik nie urósł.}
